
\subsubsection{Neural Network Calibration}

Temperature scaling emerged as the dominant post-hoc calibration method after Guo et al. (2017) demonstrated that modern neural networks are poorly calibrated despite high accuracy. Their work established Expected Calibration Error as the standard metric and showed that a single learned temperature parameter substantially reduces miscalibration. Class-based temperature scaling \citep{frenkel2021class} demonstrated that per-class temperatures can outperform global scaling in vision tasks.

\subsubsection{LLM Uncertainty Quantification}

Kadavath et al. (2022) studied whether models can identify questions they cannot answer correctly. Xiong et al. (2023) compared verbalized confidence against logit-based probabilities, finding that different elicitation methods yield different calibration properties. A key finding from this literature: LLMs trained with RLHF exhibit systematic overconfidence.

\subsubsection{Truthfulness Benchmarks}

TruthfulQA \citep{lin2022truthfulqa} provides the primary benchmark for this work. It contains 817 questions across 38 categories designed to elicit false statements that humans find plausible. No prior work has disaggregated calibration metrics by category on TruthfulQA.
